\documentclass[a4paper]{scrartcl}
\usepackage[utf8]{inputenc}
\usepackage{amsmath,amssymb,amsthm}
\usepackage{enumitem,setspace}
\usepackage[usenames,dvipsnames,table]{xcolor}
\usepackage{graphicx}
\DeclareGraphicsExtensions{.pdf,.png,.eps,.jpg}
\usepackage{algorithm2e}

% --- math operators ---
\usepackage{mathtools}
\DeclarePairedDelimiter\set{\{}{\}} % use $\set{1, 2, 3}$
\DeclarePairedDelimiter\abs{\lvert}{\rvert}
\DeclarePairedDelimiter\norm{\lVert}{\rVert}
\DeclarePairedDelimiter\ceils{\lceil}{\rceil}
\DeclarePairedDelimiter\floor{\lfloor}{\rfloor}
\def\then{\ensuremath{\Rightarrow}} % =>
\def\iff{\ensuremath{\Leftrightarrow}} % if and only if <=>
\def\to{\ensuremath{\rightarrow}} % ->
\def\ot{\ensuremath{\leftarrow}} % ->
\def\Oh{\ensuremath{\mathcal{O}}} % big O like $\Oh(n)$
% ---
\DeclareMathOperator{\preorder}{preorder}
\DeclareMathOperator{\postorder}{postorder}
\DeclareMathOperator{\depth}{depth}

% --- FILL OUT THESE 3 FIELDS ---
\def\sheetNumber{1}
\def\names{NAMEN} 
\def\sumPoints{20} 

% --- ### BEGIN DOCUMENT ### ---
\begin{document} 

\begin{doublespace} % header
\noindent \textbf{\LARGE{Übungsblatt \sheetNumber}}
\hfill  {\large Punkte: $\boxed{\qquad  /\; \sumPoints}$}\\
{\Large \textbf{Fortgeschrittene Algorithmen (Wintersemester 2022/23)}\\
Bearbeitet von: \textbf{\names}\\
\today}
\end{doublespace}


\section*{Aufgabe 1}

\section*{Aufgabe 2}
    a)
\[
\textbf{Gegeben:} \quad
\begin{cases}
\text{Kaufpreis: } M,\\[4pt]
\text{Zwei-Münzen-Strategie: }
\begin{cases}
s_1 = M/2 & \text{mit Wahrscheinlichkeit } p = 1/4,\\
s_2 = M   & \text{mit Wahrscheinlichkeit } 1 - p = 3/4.
\end{cases}
\end{cases}
\]

Für eine feste tatsächliche Skidauer \(t\) gilt für eine deterministische Schwelle \(s\):
\[
\text{Kosten}(s,t) =
\begin{cases}
t, & t \le s,\\
s + M, & t > s.
\end{cases}
\]

Damit ergibt sich der Erwartungswert der Kosten des randomisierten Algorithmus zu
\[
\mathbb{E}[\text{ALG}(t)] =
\frac{1}{4} \cdot \text{Kosten}(s_1,t) +
\frac{3}{4} \cdot \text{Kosten}(s_2,t).
\]
Das Offline-Optimum ist \(\text{OPT}(t) = \min(t, M)\).

---

\textbf{Fall 1:} \(t \le \frac{M}{2}\) \\
Beide Strategien kaufen nicht:
\[
\mathbb{E}[\text{ALG}] = t, \quad
\frac{\mathbb{E}[\text{ALG}]}{\text{OPT}} = 1.
\]

---

\textbf{Fall 2:} \(\frac{M}{2} < t \le M\) \\
Die erste Strategie kauft, die zweite nicht:
\[
\mathbb{E}[\text{ALG}] =
\frac{1}{4}(M/2 + M) + \frac{3}{4}t
= 0.375M + 0.75t,
\quad \text{OPT} = t.
\]
Somit:
\[
\frac{\mathbb{E}[\text{ALG}]}{\text{OPT}}
= 0.375 \frac{M}{t} + 0.75.
\]
Da dies für \(t\) fallend ist, liegt das Maximum bei \(t \to M/2^+\):
\[
\frac{\mathbb{E}[\text{ALG}]}{\text{OPT}}
= 0.375 \cdot 2 + 0.75 = 1.5.
\]

---

\textbf{Fall 3:} \(t > M\) \\
In beiden Fällen wird gekauft:
\[
\mathbb{E}[\text{ALG}] =
\frac{1}{4}(M/2 + M) + \frac{3}{4}(M + M)
= 0.375M + 1.5M = 1.875M,
\]
\[
\text{OPT} = M \quad \Rightarrow \quad
\frac{\mathbb{E}[\text{ALG}]}{\text{OPT}} = 1.875 = \frac{15}{8}.
\]

---

\[
\boxed{
\text{Das erwartete kompetitive Verhältnis beträgt } \frac{15}{8} = 1.875.
}
\]

\bigskip
b) Im ersten Schritt wurde auf alle Seiten im Cache 0 mal zugegriffen. Fragt man nun eine Seite an, die nicht im Cache ist, 
wird eine Seite aus dem Cache entfernt. Da alle Seiten gleich oft benutzt wurden, wird die erste Seite entfernt. Fragt 
man nun die Seite an die gerade entfernt wurde wird Seite 2 entfernt, diese am wenigsten genutzt wurde.
Duch kontinuierliches anfragen von Seiten die im vorherigen Schritt entfernt wurden, wird in jedem
Schritt ein Seitenfehler produziert.
% \section*{Aufgabe 3}

% \section*{Aufgabe 4}


% -- code for figures
% \begin{figure}[htb]
% \centering
% \includegraphics{filename} % use option [width=\linewidth] to change size to linewidth (or 0.6\linewidth for example)
% \label{fig:name}
% \end{figure}


\end{document}
% --- ### END DOCUMENT ### ---